\clearpage
\section{Copying and store: foundations for Proposition 12.2}
\label{sec:copying-foundations}
\Needspace{14\baselineskip}
\subsection{Scope, dependency map, and semantic boundary}

The new namespace \texttt{Hypernets.Copying} begins the extension needed for
Proposition~12.2 of the PDF (Proposition~12 in the HTML). The requested endpoint
is the observational equivalence of the private-cell program and the constant
program in every distinct-list environment. That endpoint is \emph{not yet
proved}. In particular, the new syntax is not an operational semantics, and
its substitution equations are not unrestricted call-by-value beta rules.

The following dependency map distinguishes checked prerequisites from pending
semantic work. The paper references use the PDF numbering, with HTML aliases
where different.
\begin{enumerate}
\item \textbf{Checked foundations:} Figure~13 syntax and substitution algebra;
term/atom wire distinction and box interfaces from Sections~9--10; concrete
canonical trees of Definition~10.1; recursive renaming, addressed replacement,
and preservation of global label properties. Ordered distributors, parallel
assembly and serial assembly are now also checked; see
Section~\ref{sec:copying-assembly}. Arbitrary contraction trees, root
termination, and uniqueness of atom-store resolution are proved in
Section~\ref{sec:copying-tracing}. All Figure~16 clauses, the full recursive translation, and global alphabet
properties are checked in Sections~\ref{sec:literal-translation}--\ref{sec:complete-translation}.
Copyable and stable graphs (Definitions~10.2--10.3), their structural
properties and recursive-isomorphism invariance are proved in
Section~\ref{sec:copying-values}.
\item \textbf{Graph correspondence:} context translation is defined and its
constructor interpretation laws are checked. Actual filling is proved for
every original annotated context, with cumulative padding, sole-hole and
binding-free certificates (Section~\ref{sec:copying-context-filling}).
Padding preserves and reflects stability and preserves copyability
(Section~\ref{sec:copying-padding-predicates}).
Distributor resource permutations are proved in
Section~\ref{sec:copying-distributor-permutation}; this is a component of the
renaming argument, not whole-term renaming.
Section~\ref{sec:copying-boundary-iso} adds a separate typed boundary-isomorphism
interface and the original distributor's explicit resource automorphism.
Selected-resource and all atomic bijective-renaming cases are proved
(Section~\ref{sec:copying-atomic-renaming}). Full-syntax bijective renaming is now proved by explicit BPerm derivations
(Section~\ref{sec:copying-full-renaming}); their exact-cost contextual
interpretation and the resulting semantic bijective-renaming law are now
proved in Section~\ref{sec:copying-semantic-renaming}. Non-bijective graph
renaming, graph-level substitution and the remaining operational invariants
remain pending. Recursive label safety and stable-input
inertness are now proved in Section~\ref{sec:copying-selection}.
\item \textbf{Checked machine laws and coordinate specimen instance:} operation search and the
extra search rules (Figure~18) now have concrete incidence rules, determinism,
typing, value recognition, exact isomorphism/padding transport and shallow
plugging preservation (Section~\ref{sec:copying-search-phase}). Reference creation (Figure~19(a)) now has its actual stable-argument pair,
freshness and resolution laws, and a search/allocation execution subsystem
(Section~\ref{sec:copying-allocation}). Dereferencing (Figure~19(b)) now has a
general local rule for arbitrary atom contraction trees, store-resolution
preservation and an execution to copy request
(Section~\ref{sec:copying-dereference}). The copy rule (Figure~17) now has a literal finite-graph construction,
local execution, safety, renaming, forward padding and shared-store resolution
(Section~\ref{sec:copying-rule}). Stable assignment (Figure~19(c)) now has a
concrete local rule, preserved store resolution and old contents, checked
execution, safety, renaming and forward padding
(Section~\ref{sec:copying-assignment}). Both equality-testing cases (Figure~19(d)) now have concrete rules,
exact local runs, safety, representation laws and context-stable store-identity
outcomes (Section~\ref{sec:copying-atom-equality}). Stable micro-beta application now has a concrete rule and quantitative local
executions (Section~\ref{sec:copying-application}). All six local machines now embed in one shallow context-closed subsystem,
with exact run costs and term-token typing. Arithmetic rules now extend the
same relation, with five-step unary and seven-step binary executions. Its concrete allowance-indexed
execution has exact finite-run correspondence and proved recursive
isomorphism invariance (Section~\ref{sec:copying-combined}). Every combined run
and closed observation pads forward at exactly the same cost
(Section~\ref{sec:copying-arithmetic-padding}). Every padded match and run now
also reflects at exactly the same cost; copyability reflection and preservation
of representability are proved from the concrete graph rules
(Section~\ref{sec:copying-padding-reflection}). Canonical rule/operand selection
is proved for independently chosen matches in arbitrary hosts; so are copy-head
and read/write-store identity, disjoint equality branches, and recursive
runtime safety for complete finite runs (Section~\ref{sec:copying-selection}).
Allocation now has complete successor-state coherence for independently
chosen contextual matches, with exact-cost consequences and a unique successful
cost for the literal translated reference expression
(Section~\ref{sec:copying-allocation-coherence}). Both reference-equality
branches now have complete successor-state coherence, including arbitrary
hosts, whole store attachments and token positions. Their full-machine
successful cost is exactly seven with arbitrary contents and finite padding
(Section~\ref{sec:copying-equality-coherence}). Assignment now has complete
successor coherence too, retaining the store, updated payload, discarded old
contents and token position. Typed target exchange commutes with arbitrary
context formation; stable writes have exactly seven successful steps with
arbitrary attached contents and translated-data padding regressions
(Section~\ref{sec:copying-assignment-coherence}). Dereference now also has
complete contextual successor coherence, closing this obligation for all
store-operation heads. Its exact cost law retains the five-step read prefix
and the subsequent copying continuation, including unchanged translated data
at every finite padding (Section~\ref{sec:copying-dereference-coherence}).
Arithmetic now has complete contextual successor coherence, exact five- and
seven-step successful costs, and an exact seven-step handoff to an arbitrary
host continuation (Section~\ref{sec:copying-arithmetic-coherence}).
Application has complete contextual successor coherence, including
independently represented bodies and arbitrary contexts, with an exact six-step
handoff (Section~\ref{sec:copying-application-coherence}). Copying now completes
the final successor comparison, including independent tree boundaries, complete
recursive payloads, original distributors and retained hosts. The full finite
machine has successor determinism up to recursive state isomorphism, and the
unbounded machine up to its representation relation. Exact prefix costs and
unique successful-run lengths are proved without extra machine-law premises
(Section~\ref{sec:copying-successor-coherence}).

Section~\ref{sec:copying-intrinsic-refocusing} proves refocusing for search,
copying and allocation, together with exact search history, single-head
replacement and rootedness representation/padding laws.
Section~\ref{sec:copying-full-refocusing} now discharges all remaining
operation cases: dereference, assignment, reference equality, application and
arithmetic. The full finite and unbounded machines are refocusing without any
compute-law premise. Every reachable state, including every finite prefix of
an infinite execution, has a pure-search witness on its own current graph.
Section~\ref{sec:copying-specimens} constructs the coordinate specimen system
and proves synchronized-search locality for the complete unbounded machine.
Recursive paired fillings retain full box bodies and ordered interfaces.
The machine is literally deterministic on classes under recursive isomorphism
and depth padding; every class run reflects from any fixed original source
with its exact length. All existing binding-free observations are unchanged.
Section~\ref{sec:copying-contextual-robustness} now constructs the recursive
certificates for every original addressed closing-context filling, including
holes inside boxes. Their actual initial states supply the required lifting,
so both quantitative contextual robustness conclusions cover the entire
original observation class. Term observations, cofinal allowances and the
full term-context grammar are connected to this theorem using the unchanged
translation. Template output closure, safety and robustness remain application
proofs. The remaining correspondence work concerns explicitly focused graphs,
unique-hole inventories and replacement tuples. The ten other template families,
semantic weakening and final Figure~22 argument remain open.
Section~\ref{sec:copying-box-permutation-opening} discharges BPerm output
closure and operational input safety, and proves its local beta-opening diagram
with equal transition counts and recursively isomorphic successors. External
interface permutations now preserve exact successful costs for the entire
combined machine, through the original context-closure rule. Full BPerm robustness is completed in Section~\ref{sec:copying-full-bperm}. Section~\ref{sec:copying-box-permutation-overlaps}
now proves the local shallow copying diagram, including both surviving boxes
and the original resource distributors, and strengthens beta opening to equal
exact costs in arbitrary typed hosts from the redex input.
Section~\ref{sec:copying-residual-specimens} now constructs genuine recursive
residual certificates for shallow and all-deep copying, composes mixed
contributions, and packages closed residuals as original quasi-specimens.
The closed all-deep copying case has forward/converse robustness conclusions.
Section~\ref{sec:copying-beta-residuals} proves the local deep beta transfer,
including padding of genuine certificates, preservation of stability, the
stronger rigidity of stable BPerm fillings, and both closed robustness cases.
A shallow box can be paired with an independently certified stable argument,
and a mixed-host regression retains a separate nontrivial contribution.
Section~\ref{sec:copying-overlap-classification} now exhaustively classifies
actual rewrite sources in arbitrary original specimens. Primitive footprints
have constructed opposite matches; beta frames retain their original box
certificates, and sharing paths preserve their exact lengths. Micro-beta and
reference creation also have proved output-closed pre-templates tied to the
original compute rules. Section~\ref{sec:copying-full-bperm} now completes
BPerm matching whole-graph rewrites and original pointed residual certificates
for every classified case, and proves full forward/converse robustness.
The exact-cost consequence holds for arbitrary original closed paired fillings;
none of the remaining template obligations is assumed in this proof.
The concrete instance now supplies every abstract machine-law and specimen
locality field. Output closure, safety and robustness remain conditions to
prove for each template. Copy transitions retain their own concrete rules;
they are not hidden in the linear machine's existing compute parameter.
\item \textbf{Contextual interface:} binding-free closure at sourced holes,
binding-freedom of all translated term contexts, sole-hole certificates and
addressed context composition are checked. Graph-context congruence has an
explicit execution-isomorphism premise, now discharged for the combined
subsystem, including arithmetic. Observation reflection and the cofinal
allowance criterion from Section~\ref{sec:copying-observation-padding} now have
both execution premises discharged. Section~\ref{sec:copying-padding-reflection}
instantiates full term-context congruence with the concrete execution, without
an equal-minimum-depth restriction. These representation laws must be preserved
through any semantic correction required by the remaining full-machine and
paper-correspondence obligations of Section~\ref{sec:paper-correspondence}. Unrestricted
zero-input-hole closure is refuted, as recorded in Appendix~B.1.
\item \textbf{Completed template obligations:} BPerm output closure, input safety for arbitrary cost parameters, and its local beta-opening diagram; see Section~\ref{sec:copying-box-permutation-opening}. External compensation has exact machine-cost laws. The local shallow copying
diagram and exact whole-host beta-redex costs are also checked
(Section~\ref{sec:copying-box-permutation-overlaps}). Recursive copying residual
certificates, mixed certificate assembly, framed residual transfer and the
closed all-deep copying robustness cases are checked in
Section~\ref{sec:copying-residual-specimens}. Padding of actual certificates,
stable-filling rigidity, local deep beta residuals, closed forward/converse
beta robustness cases and shallow beta with a separately certified argument
are checked in Section~\ref{sec:copying-beta-residuals}. Source-overlap
classification for arbitrary actual rewrites, primitive occurrence reconstruction,
beta-frame and counted-path provenance, and output closure of micro-beta and
reference creation are proved in Section~\ref{sec:copying-overlap-classification}.
Full BPerm whole-graph matching, general pointed residual assembly and both
robustness directions are proved in Section~\ref{sec:copying-full-bperm},
with exact successful-transition counts for closed paired fillings.
Micro-beta and reference creation now also have input safety for unrestricted
and binding-free contexts, retaining the six/four-step versus zero-step
comparison (Section~\ref{sec:copying-operational-safety}).
Section~\ref{sec:copying-operational-transfer} adds constructive finite
complements and whole-graph rewrite factorization, retained matched
skeletons with actual recursive certificates, counted copying barriers,
and framed deep beta/copying forward and converse robustness with
binding-free residuals. Padding preserves the actual shared context by a
proved recursive isomorphism.
\item \textbf{Pending templates:} Table~3 requires contraction associativity,
commutativity and idempotence; copying and garbage collection;
pulling contractions and pulling weakening; micro-beta; reference creation;
and parametricity. Micro-beta and reference creation have output closure and
input safety; their forward/converse robustness remains pending. The other
eight families still require their template obligations. These proofs use the transfer results
of Appendices~G--H. Prove base structural templates before the dependent
copying, box-pulling and parametricity arguments; do not assume their contextual
consequences in a circular proof.
\item \textbf{Pending assembly:} the semantic weakening laws (Proposition~12.1;
HTML Proposition~11), the empty-environment chain of Figure~22, and extension
to arbitrary environments (Proposition~12.2; HTML Proposition~12). The copying
and garbage-collection templates are indirect dependencies, as Table~4 records.
\end{enumerate}
The final public theorem is reserved as
\texttt{Copying.Examples.proposition12\_2}; no placeholder theorem with that name
is introduced. Definitions and theorems below are reusable prerequisites rather
than a conditional assertion of the requested conclusion.

\Needspace{14\baselineskip}
\subsection{Scoped terms and the two programs}

\termdef{copying-scoped-term}
Constants are naturals, booleans and unit. Unary syntax contains numeric
negation, allocation and dereference; binary syntax contains addition,
subtraction, assignment and atom equality. A term $t\in\mathsf{Term}(k,h)$ has
$k$ available variables and $h$ available atoms. Variables and atoms are finite
indices in separate namespaces. A lambda body has $k+1$ variable slots; the
newly bound variable has index zero. Both children of application and binary
operations use the same environments, so variable sharing is permitted.
These are formation rules, without an evaluation relation.

\lean{Copying/Syntax}{7}{11}{Syntax.Constant}

\lean{Copying/Syntax}{13}{15}{Syntax.Unary}

\lean{Copying/Syntax}{17}{20}{Syntax.Binary}

\lean{Copying/Syntax}{23}{32}{Syntax.Term}

Lifting a renaming fixes index zero and maps each successor index by the old
renaming followed by successor. The zero and successor equations hold by
definition. Simultaneous renaming recursively acts on variables and atoms;
only the variable map is lifted under lambda. Non-injective maps are allowed,
as identifying variable names is a legitimate substitution operation here.

\lean{Copying/Syntax}{34}{35}{Syntax.lift}

\lean{Copying/Syntax}{37}{37}{Syntax.lift_zero}

\lean{Copying/Syntax}{38}{40}{Syntax.lift_succ}

\lean{Copying/Syntax}{42}{49}{Syntax.rename}

Lifting preserves identity and composition: split a finite index into zero
or successor and both statements reduce to equality by definition. Structural
induction then proves identity and composition for renaming all terms. The
lambda cases use the corresponding lifting identities; every other constructor
uses the induction hypotheses on its children.

\lean{Copying/Syntax}{51}{52}{Syntax.lift_id}

\lean{Copying/Syntax}{54}{56}{Syntax.lift_comp}

\lean{Copying/Syntax}{58}{61}{Syntax.rename_id}

\lean{Copying/Syntax}{63}{70}{Syntax.rename_comp}

Weakening inserts a fresh variable at the front, shifting every old variable
index by one and leaving atoms unchanged. The let form is application of a
lambda to its bound expression. Sequencing $u;t$ is this let form with weakened
$t$, ensuring the fresh sequencing binder does not capture any variable in $t$.

\lean{Copying/Syntax}{72}{73}{Syntax.weaken}

\lean{Copying/Syntax}{75}{76}{Syntax.letIn}

\lean{Copying/Syntax}{78}{79}{Syntax.seq}

A lifted substitution maps zero to the new bound variable and maps a successor
to the weakened old replacement. Substitution recursively replaces variable
leaves with terms, renames atom leaves, and lifts the replacement environment
under lambda. Replacements are not substituted into again by this same
operation. This is capture-avoiding simultaneous syntactic substitution.

\lean{Copying/Syntax}{81}{83}{Syntax.liftSub}

\lean{Copying/Syntax}{85}{92}{Syntax.subst}

Lifting a substitution consisting only of variable leaves equals lifting the
underlying renaming. Splitting the index proves this equation. Induction on
terms then identifies substitution by variables with renaming. Its identity
instance follows from the previously proved renaming identity.

\lean{Copying/Syntax}{94}{96}{Syntax.liftSub_var}

\lean{Copying/Syntax}{98}{102}{Syntax.subst_var}

\lean{Copying/Syntax}{104}{106}{Syntax.subst_id}

The definitions \texttt{privateCell} and \texttt{constantResult} encode the two
programs in every scope $(k,h)$. In the desugared left program, the dereferenced
$x$ has index two beneath the binders for $f$ and the sequencing variable.
The two desugaring theorems are definitional equalities. Neither program uses
an ambient variable or atom: applying any ambient renaming produces exactly
the same program in the target scope. These equalities also hold by reduction
of the syntax and lifting definitions, not by an observational argument.

\lean{Copying/Syntax}{108}{111}{Syntax.privateCell}

\lean{Copying/Syntax}{113}{115}{Syntax.constantResult}

\lean{Copying/Syntax}{117}{119}{Syntax.privateCell_desugar}

\lean{Copying/Syntax}{121}{123}{Syntax.constantResult_desugar}

\lean{Copying/Syntax}{125}{126}{Syntax.privateCell_rename}

\lean{Copying/Syntax}{128}{131}{Syntax.constantResult_rename}

A named environment consists of ordered variable and atom lists with separate
no-duplicates proofs. The two lists need not be disjoint: they designate separate
namespaces. Their lengths determine the indices of the scoped term type.
This retains the distinct-list presentation of the requested statement without
requiring fresh printed binder names inside the implementation.

\lean{Copying/Syntax}{134}{138}{Syntax.Environment}

\Needspace{14\baselineskip}
\subsection{Substitution algebra}

Renaming a weakened term under a lifted renaming equals weakening its renamed
term. Expand weakening, apply renaming composition, and compare the two maps
pointwise. Splitting zero/successor lifts this identity to substitution
environments. Structural induction now shows that renaming after substitution
renames every replacement and composes the atom maps. The lambda case is
exactly the lifted-environment identity.

\lean{Copying/Substitution}{6}{9}{Substitution.rename_weaken}

\lean{Copying/Substitution}{11}{18}{Substitution.rename_liftSub}

\lean{Copying/Substitution}{20}{26}{Substitution.rename_subst}

Reindexing a lifted replacement family equals lifting its reindexed family;
zero and successor prove the assertion directly. Consequently substitution
after renaming simply reindexes the replacement family and composes atom maps.
The proof again follows term induction, using the family equation under lambda.

\lean{Copying/Substitution}{28}{31}{Substitution.liftSub_rename}

\lean{Copying/Substitution}{33}{39}{Substitution.subst_rename}

Substitution under a fresh binder commutes with weakening. This follows by
applying the two preceding interchange laws to the successor renaming. Splitting
an index lifts the result to replacement families. Finally, two successive
substitutions equal one substitution whose replacements are themselves
substituted, and whose atom map is the composite. Induction on the original
term proves this associativity law; the lambda case uses the lifted-family law.
Thus the binder cases have been discharged, rather than left as freshness
assumptions.

\lean{Copying/Substitution}{41}{44}{Substitution.subst_weaken}

\lean{Copying/Substitution}{46}{53}{Substitution.subst_liftSub}

\lean{Copying/Substitution}{55}{62}{Substitution.subst_comp}

Instantiation replaces variable zero with a supplied term and keeps older
variables, shifting them down by one. Instantiating a weakened term returns
the original term, by substitution--renaming interchange and substitution
identity. Instantiating variable zero returns the supplied term by definition.
These syntactic identities do not license reducing a call-by-value application
before its argument is a value.

\lean{Copying/Substitution}{64}{65}{Substitution.instantiate}

\lean{Copying/Substitution}{67}{70}{Substitution.instantiate_weaken}

\lean{Copying/Substitution}{73}{73}{Substitution.instantiate_zero}

\Needspace{14\baselineskip}
\subsection{Copying signatures and finite boxed graphs}

\termdef{copying-wire}
Wire types are term, atom, and a thunk annotated by its number of bound
variables. Sharing sorts are only term and atom; the inclusion into wire types
therefore cannot select a thunk. Operations comprise all constants, lambda,
application, and all unary/binary syntax operations. Constants and lambda are
passive; the remaining operations are active. Every operation has one term
input. Outputs are empty for constants, a one-binder thunk for lambda, two
terms for application/binary operations, and one term for unary operations.

\lean{Copying/Hierarchy}{10}{13}{Hierarchy.Wire}

\lean{Copying/Hierarchy}{15}{17}{Hierarchy.ShareSort}

\lean{Copying/Hierarchy}{19}{22}{Hierarchy.ShareSort.wire}

\lean{Copying/Hierarchy}{24}{29}{Hierarchy.Operation}

\lean{Copying/Hierarchy}{31}{33}{Hierarchy.Operation.kind}

\lean{Copying/Hierarchy}{35}{39}{Hierarchy.Operation.outputs}

Holes have names and ordered typed interfaces. Primitive edges are operations,
contractions, weakenings, instances, stores, focuses and holes. Contraction
has two inputs and one output of its sharing sort; weakening has no inputs
and one output. An instance has type $\star\Rightarrow\diamond$, and store has
the opposite type. Focus has type $\star\Rightarrow\star$. The distinction
between term and atom contractions is present in the labels; no copy rule for
atom contractions is supplied. There is no primitive linear variable edge.

\lean{Copying/Hierarchy}{41}{46}{Hierarchy.Hole}

\lean{Copying/Hierarchy}{48}{55}{Hierarchy.Primitive}

\lean{Copying/Hierarchy}{57}{70}{Hierarchy.primitiveSignature}

The environment interface lists all term outputs before all atom outputs.
An alphabet is a label type with its signature. A box records a bound-variable
count, a free-variable count, a free-atom count, and a body whose outputs are
bound variables, free variables, then free atoms. Its external input is the
appropriate thunk, and its external outputs omit the bound variables.
The zero-depth alphabet contains primitives; a successor alphabet adds boxes
whose bodies use the previous alphabet. A hypernet is the existing finite
interfaced graph over this alphabet, retaining every incidence and typing law.

\lean{Copying/Hierarchy}{72}{73}{Hierarchy.environmentInterface}

\lean{Copying/Hierarchy}{75}{78}{Hierarchy.Alphabet}

\lean{Copying/Hierarchy}{80}{84}{Hierarchy.Box}

\lean{Copying/Hierarchy}{86}{91}{Hierarchy.baseAlphabet}

\lean{Copying/Hierarchy}{93}{102}{Hierarchy.nextAlphabet}

\lean{Copying/Hierarchy}{104}{106}{Hierarchy.alphabet}

\lean{Copying/Hierarchy}{108}{108}{Hierarchy.Hypernet}

Primitive labels embed into every finite alphabet. Both signature equations
are proved by splitting the depth into zero or successor. The primitive graph
is the existing one-edge corolla at the embedded label; those equations give
its stated interface without changing its incidence.

\lean{Copying/Hierarchy}{110}{112}{Hierarchy.primitiveLabel}

\lean{Copying/Hierarchy}{114}{116}{Hierarchy.primitive_input}

\lean{Copying/Hierarchy}{118}{121}{Hierarchy.primitive_output}

\lean{Copying/Hierarchy}{123}{125}{Hierarchy.primitiveNet}

The operation predicate recognises only primitive operation labels. Their
sources have length one, by graph typing and case analysis on labels. Hence
the existing generic operation-path uniqueness theorem applies to copying
hypernets as well: two operation paths to the same vertex from inputs have
the same edges and the same starting input. A contraction is explicitly not
an operation, so its two sources do not violate this argument.

\lean{Copying/Hierarchy}{127}{131}{Hierarchy.IsOperation}

\lean{Copying/Hierarchy}{134}{144}{Hierarchy.unary_operations}

\lean{Copying/Hierarchy}{146}{151}{Hierarchy.operation_path_unique}

\lean{Copying/Hierarchy}{153}{155}{Hierarchy.contraction_not_operation}

Boxing constructs a one-edge graph with the given typed body. Raising increases
the nesting allowance, relabelling primitives and recursively raising box
bodies while retaining all three counts. It changes neither outer carriers
nor incidence. Heterogeneous equalities of source and target functions follow
by splitting the old depth; the heterogeneous formulation accommodates their
dependent carrier types.

\lean{Copying/Hierarchy}{157}{162}{Hierarchy.box}

\lean{Copying/Hierarchy}{164}{171}{Hierarchy.raise}

\lean{Copying/Hierarchy}{173}{174}{Hierarchy.raise_source}

\lean{Copying/Hierarchy}{176}{178}{Hierarchy.raise_target}

Deep edge occurrences are either shallow edges or a pair consisting of an
outer box edge and a child occurrence. Equal box labels at different edges
therefore still give different occurrences. Sum and dependent-sum enumeration
construct an explicit finite enumeration recursively. Every occurrence belongs
to it by the enumeration's completeness field.

\lean{Copying/Hierarchy}{180}{185}{Hierarchy.AllEdge}

\lean{Copying/Hierarchy}{187}{193}{Hierarchy.allEdgesFinite}

\lean{Copying/Hierarchy}{195}{196}{Hierarchy.every_edge_enumerated}

\Needspace{14\baselineskip}
\subsection{Recursive isomorphism and addressed replacement}

Box relatedness requires identical bound, free-variable and free-atom counts,
and related bodies at that exact interface. Label equivalence is literal on
primitives and recursive graph isomorphism on box bodies. Mixed primitive/box
labels are unrelated. A deep isomorphism uses the generic incidence-isomorphism
structure with this label relation; it preserves ordered boundaries at every
level while permitting independent internal renaming.

\lean{Copying/RecursiveIso}{9}{13}{RecursiveIso.BoxRelated}

\lean{Copying/RecursiveIso}{16}{23}{RecursiveIso.LabelEquiv}

\lean{Copying/RecursiveIso}{26}{27}{RecursiveIso.DeepIso}

Induction on depth proves reflexivity, symmetry and transitivity of label
equivalence. Primitive cases are the corresponding equality laws. In a box
case, use the induction hypothesis to build, reverse or compose the body's
incidence isomorphism. Incompatible label cases are impossible.

\lean{Copying/RecursiveIso}{29}{39}{RecursiveIso.labelEquiv_refl}

\lean{Copying/RecursiveIso}{41}{57}{RecursiveIso.labelEquiv_symm}

\lean{Copying/RecursiveIso}{59}{84}{RecursiveIso.labelEquiv_trans}

The graph-level identity, inverse and composition apply the generic
isomorphism constructors with the three label laws just proved. An ordinary
literal-label isomorphism embeds into a deep isomorphism using reflexivity
of label equivalence.

\lean{Copying/RecursiveIso}{86}{86}{RecursiveIso.refl}

\lean{Copying/RecursiveIso}{88}{88}{RecursiveIso.symm}

\lean{Copying/RecursiveIso}{90}{92}{RecursiveIso.trans}

\lean{Copying/RecursiveIso}{94}{97}{RecursiveIso.ofIso}

Related labels have equal input and output signatures. For primitives this
follows by applying the signature function to label equality. For boxes the
relatedness constructor identifies all interface counts. The result follows
by cases, with mixed labels eliminated by contradiction.

\lean{Copying/RecursiveIso}{99}{112}{RecursiveIso.labelEquiv_input}

\lean{Copying/RecursiveIso}{114}{126}{RecursiveIso.labelEquiv_output}

Replacement of a selected box changes just its body, preserving its metadata
and all outer incidence. A typed site is either a shallow selected edge or a
site recursively inside a selected box. Its indices carry the replacement's
local allowance and interface. Addressed plugging performs the existing finite
graph substitution at a shallow site and reconstructs enclosing boxes at an
inner site. It is defined for arbitrary selected edges, not only holes.

\lean{Copying/Plugging}{11}{30}{Plugging.replaceBox}

\lean{Copying/Plugging}{33}{42}{Plugging.Site}

\lean{Copying/Plugging}{45}{51}{Plugging.plugAt}

Shallow addressed plugging equals ordinary graph plugging by definition.
Inside a box, source and target functions of the outer graph remain unchanged,
also by definition. Every other outer label is unchanged: unfold replacement
and use the inequality of the selected and retained edge. Equal contents of
distinct boxes do not trigger simultaneous replacement.

\lean{Copying/Plugging}{53}{58}{Plugging.plugAt_shallow}

\lean{Copying/Plugging}{60}{64}{Plugging.plugAt_inside_source}

\lean{Copying/Plugging}{66}{71}{Plugging.plugAt_inside_target}

\lean{Copying/Plugging}{74}{79}{Plugging.plugAt_inside_other_label}

\Needspace{14\baselineskip}
\subsection{Global predicates and closure}

\termdef{copying-global-label}
A predicate on primitive labels extends recursively through every box body.
The graph predicate requires it on all outer labels and therefore at every
depth. Focus freedom excludes all focus modes, and hole freedom excludes all
hole labels. These are global predicates, not shallow tests.

\lean{Copying/Predicates}{10}{13}{Predicates.LabelAll}

\lean{Copying/Predicates}{15}{16}{Predicates.AllLabels}

\lean{Copying/Predicates}{18}{19}{Predicates.FocusFree}

\lean{Copying/Predicates}{21}{23}{Predicates.HoleFree}

\termdef{copying-replacement-closure}
A graph class is closed under replacement when replacing any well-typed
addressed site of one member by another member produces a member. This is
stronger than closure only at holes. Global label predicates satisfy it:
at a shallow site every resulting edge comes from either the host or the
replacement; inside a box apply the induction hypothesis to the selected body
and leave other labels unchanged. Focus-free and hole-free closure are the
two immediate instances. No binding-free path condition is asserted here.

\lean{Copying/Predicates}{26}{29}{Predicates.ClosedUnderReplacement}

\lean{Copying/Predicates}{31}{51}{Predicates.allLabels_plugAt}

\lean{Copying/Predicates}{53}{54}{Predicates.focusFree_closed}

\lean{Copying/Predicates}{56}{58}{Predicates.holeFree_closed}

Raising preserves and reflects a global label property. Induction on depth
reduces the primitive case to identity and the box case to the induction
hypothesis. Recursive label equivalence likewise preserves and reflects the
property: map edges through the body's bijection, using its inverse for the
reverse coverage direction. Applying this label fact to the outer edge
bijection proves invariance under deep graph isomorphism.

\lean{Copying/Predicates}{60}{76}{Predicates.allLabels_raise_iff}

\lean{Copying/Predicates}{78}{99}{Predicates.labelAll_equiv}

\lean{Copying/Predicates}{101}{108}{Predicates.allLabels_iso_iff}

\Needspace{14\baselineskip}
\subsection{The ordered canonical contraction tree}

\termdef{copying-canonical-tree}
For $m$ inputs, internal vertices are $0,\ldots,m-1$ and edges are
$0,\ldots,m$. Edge zero is weakening. If $m=0$, it targets the output.
Otherwise edge one is the root contraction, fed on the left by the weakening.
The other contractions form the ordered right spine from Definition~10.1.
For $e>1$, its left source is input $e-2$; its right source is the next
internal vertex unless it is the last contraction, whose right source is
input $m-1$. Thus the one-input case still contains both edges. These ordered
sources follow the PDF figure, rather than merely choosing an arbitrary
binary tree with the same leaves.

\lean{Copying/CanonicalTree}{10}{12}{CanonicalTree.Vertex}

\lean{Copying/CanonicalTree}{14}{18}{CanonicalTree.node}

\lean{Copying/CanonicalTree}{20}{23}{CanonicalTree.sources}

\lean{Copying/CanonicalTree}{25}{26}{CanonicalTree.label}

Distinct edges have distinct target nodes: inspect the empty tree, weakening,
root and remaining spine cases. No target is an input. The two sources of a
contraction are distinct, and a vertex cannot be sourced by two different
edges; the proofs inspect the ordered source cases and compare finite indices.
Every non-output vertex occurs as a source: choose the root for the weakening
vertex, the matching spine edge for other internal vertices, and the appropriate
contraction for each input. Every non-input vertex is the target of its indexed
edge. Finally all vertices have the sharing sort's label, including boundary
vertices, by inspection of the repeated interface lists.

\lean{Copying/CanonicalTree}{28}{34}{CanonicalTree.node_injective}

\lean{Copying/CanonicalTree}{36}{39}{CanonicalTree.node_not_input}

\lean{Copying/CanonicalTree}{41}{46}{CanonicalTree.sources_nodup}

\lean{Copying/CanonicalTree}{48}{60}{CanonicalTree.source_unique}

\lean{Copying/CanonicalTree}{62}{89}{CanonicalTree.source_coverage}

\lean{Copying/CanonicalTree}{91}{110}{CanonicalTree.target_coverage}

\lean{Copying/CanonicalTree}{112}{116}{CanonicalTree.vertex_label}

The tree graph packages these finite carriers and incidence functions. The
preceding results discharge uniqueness and coverage. Source and target typing
follow because every vertex has the same sharing label and the primitive
signature has the corresponding arity. The label characterisations prove
that edge zero is the unique weakening and all remaining $m$ edges are
contractions.

\lean{Copying/CanonicalTree}{118}{142}{CanonicalTree.net}

\lean{Copying/CanonicalTree}{144}{150}{CanonicalTree.weakening_iff}

\lean{Copying/CanonicalTree}{152}{159}{CanonicalTree.contraction_iff}

Two tree edges are adjacent exactly when they are weakening followed by the
root, or consecutive edges in the upward right spine. Unfold singleton target
membership and the source cases to prove this equivalence. Assign weakening
rank zero and a nonzero edge $e$ rank $m+1-e$. Each adjacency strictly increases
rank, by the just-proved characterisation and arithmetic.

\lean{Copying/CanonicalTree}{161}{169}{CanonicalTree.adjacent_iff}

\lean{Copying/CanonicalTree}{171}{171}{CanonicalTree.rank}

\lean{Copying/CanonicalTree}{173}{178}{CanonicalTree.adjacent_rank_lt}

Induction on the path list proves that earlier positions have strictly smaller
rank than later positions. At its head, combine the first adjacency with the
induction hypothesis on the tail; two tail positions reduce directly to that
hypothesis. Equal edges at different positions would then have strictly
ordered equal ranks, a contradiction. Hence every path has distinct edges.
This acyclicity is a theorem for these trees, not an axiom imposed on general
store graphs. Finally signature-preserving relabelling embeds the tree into
any finite copying alphabet.

\lean{Copying/CanonicalTree}{180}{205}{CanonicalTree.path_rank_lt}

\lean{Copying/CanonicalTree}{207}{221}{CanonicalTree.path_nodup}

\lean{Copying/CanonicalTree}{223}{225}{CanonicalTree.atDepth}

\Needspace{14\baselineskip}
\subsection{Regression examples and what they establish}

The first regression substitutes an open variable beneath a lambda and checks
that its index is shifted rather than captured. The second inserts an atom
beneath a lambda and checks that the atom index is not shifted. Both equations
are proved by evaluating the recursive syntax definitions.

\lean{Copying/Examples}{11}{14}{Examples.substitution_avoids_capture}

\lean{Copying/Examples}{16}{19}{Examples.atom_not_captured}

The one-input regression checks the ordered sources: weakening first, input
second. The two-input regression checks the lower contraction's ordered input
ports. The three-input regression checks weakening-to-root and both spine
adjacencies, and excludes the reverse edge. The first two reduce directly;
the last uses the adjacency characterisation and finite arithmetic.

\lean{Copying/Examples}{21}{24}{Examples.unary_tree_sources}

\lean{Copying/Examples}{26}{29}{Examples.binary_tree_sources}

\lean{Copying/Examples}{31}{38}{Examples.ternary_tree_adjacency}

The boxed-instance graph has no bound or free term variables and one free
atom. Its interface is a zero-binder thunk to an atom output. Its body is an
instance corolla. Inspection of the sole primitive inside the box proves both
global focus freedom and global hole freedom.

\lean{Copying/Examples}{40}{41}{Examples.boxedInstance}

\lean{Copying/Examples}{43}{46}{Examples.boxedInstance_focusFree}

\lean{Copying/Examples}{48}{52}{Examples.boxedInstance_holeFree}

For every named distinct-list environment, both displayed programs are the
renamings of their empty-environment syntax. Apply the two program renaming
theorems to the unique maps from empty index sets. This does not prove the
semantic weakening laws, which compare translated graphs and remain on the
dependency map. The new axiom audit checks these prerequisites separately;
there is no assertion of Proposition~12.2 in this snapshot.

\lean{Copying/Examples}{55}{60}{Examples.programs_in_environment}


\paragraph{Binding-free closure audit.}
Section~\ref{sec:copying-binding} completes positive-path correspondence,
recursive-isomorphism invariance and plugging closure at sourced holes.
All translated term holes satisfy the source condition. The checked
zero-input counterexample qualifies unrestricted Lemma~E.1; copying
observations and lambda-context congruence are still pending.


The remaining translation dependency for context composition is discharged in
Section~\ref{sec:copying-translated-composition}: independently translated hosts
compose by actual addressed plugging, with the required padding, and both the
composite host and all its fillings agree with the original direct translation.
Graph-level term renaming and substitution remain separate obligations.


The serial-renaming dependency is discharged in
Section~\ref{sec:copying-serial-renaming}: typed boundary maps glue when their
shared actions agree, with a converse and concrete distributor instances.
Structural stability and nested binding-freedom transport are proved. The
composite-term induction, including the box-body boundary correspondence, is
still required; this milestone does not discharge concrete execution laws.


Unary translation renaming is now discharged by
Section~\ref{sec:copying-unary-renaming}, including the mandatory single-use
forests and all three unary operations. The remaining bijective cases are
application, binary operations, and lambda boxes, not unary operations.


The remaining application and binary renaming dependencies are discharged in
Section~\ref{sec:copying-lambda-free-renaming}, including exact resource matching,
the original shuffle, and independent padding to a common allowance. Every
lambda-free term now has a public graph-renaming theorem. Lambda-box interface
renaming and graph-level substitution remain separate obligations.


The lambda dependency now factors through the actual BPerm graph pair in
Section~\ref{sec:copying-full-renaming}. Its exact bound/free interface laws and
abstraction-shell compatibility complete the full syntax induction. The remaining
edge in the dependency map is semantic: interpret the explicit derivation using
BPerm robustness and execution congruence. Non-bijective graph substitution and
the copying-machine obligations remain open.
